\documentclass[10pt, a4paper]{article}

% --- Pacotes Fundamentais ---
\usepackage[utf8]{inputenc}
\usepackage[T1]{fontenc}
\usepackage[brazil]{babel}
\usepackage{geometry}
\usepackage{titlesec}
\usepackage{fancyhdr}
\usepackage{xcolor}
\usepackage{enumitem}
\usepackage{tabularx}
\usepackage{colortbl}
\usepackage{lmodern}
\usepackage{lastpage}
\usepackage{tikz}
\usepackage{parskip}
\usepackage{hyperref}
\usepackage{multicol}

% --- Configurações de Layout (margens compactas para caber em 1 página) ---
\geometry{left=1.5cm, right=1.5cm, top=2.8cm, bottom=1.5cm, headsep=0.8cm, headheight=40pt}

% --- Definição de Cores (mesma paleta do relatório de tráfego) ---
\definecolor{primaryColor}{RGB}{0, 51, 102}
\definecolor{secondaryText}{RGB}{80, 80, 80}
\definecolor{accentColor}{RGB}{178, 34, 34}
\definecolor{lightGray}{RGB}{245, 245, 245}
\definecolor{tableHeader}{RGB}{230, 230, 250}
\definecolor{successGreen}{RGB}{34, 139, 34}
\definecolor{warningOrange}{RGB}{255, 140, 0}

% --- Hyperlinks ---
\hypersetup{
    colorlinks=true,
    linkcolor=primaryColor,
    urlcolor=accentColor,
    pdftitle={Guia de Preenchimento NFS-e -- Serviços Odontológicos},
    pdfauthor={Fred Lisboa}
}

% --- Cabeçalho e Rodapé ---
\pagestyle{fancy}
\fancyhf{}

\fancyhead[C]{%
    \begin{tabularx}{\textwidth}{@{} X r @{}}
        \textcolor{primaryColor}{\textbf{\footnotesize GUIA DE PREENCHIMENTO -- NFS-e}} & \scriptsize \textit{Serviços Odontológicos} \\
        \textcolor{secondaryText}{\scriptsize Emissão de Nota Fiscal de Serviço Eletrônica} & {\scriptsize \textbf{Referência: Abril/2026}}
    \end{tabularx}%
}

\fancyfoot[L]{\scriptsize \textcolor{secondaryText}{NFS-e | \textbf{ME/EPP -- Simples Nacional}}}
\fancyfoot[R]{\normalsize Página \textbf{\thepage} de \textbf{\pageref{LastPage}}}

\renewcommand{\headrulewidth}{1.5pt}
\renewcommand{\headrule}{%
    \vspace{0.05cm}
    \hbox to\headwidth{\color{primaryColor}\leaders\hrule height \headrulewidth\hfill}%
}
\renewcommand{\footrulewidth}{1.5pt}
\renewcommand{\footrule}{%
    \color{primaryColor}%
    \hbox to\headwidth{\leaders\hrule height \footrulewidth\hfill}%
    \vspace{0.1cm}
}

% --- Formatação de Títulos ---
\titleformat{\section}
{\color{primaryColor}\normalfont\normalsize\bfseries}
{\thesection}{0.6em}{}[{\titlerule[0.6pt]}]
\titlespacing*{\section}{0pt}{6pt}{3pt}

\titleformat{\subsection}
{\color{primaryColor}\normalfont\small\bfseries}
{}{0em}{}
\titlespacing*{\subsection}{0pt}{4pt}{2pt}

% --- Espaçamento compacto ---
\setlength{\parskip}{1pt}
\setlist[itemize]{nosep, leftmargin=1.2em, itemsep=0pt, parsep=0pt, topsep=1pt}

% --- INÍCIO DO DOCUMENTO ---
\begin{document}

% --- Seção 1: Campo Descrição ---
\section{Localização do Campo ``Descrição dos Serviços''}

\colorbox{warningOrange!15}{%
    \parbox{\dimexpr\textwidth-2\fboxsep}{%
        \small\textbf{\textcolor{accentColor}{Campo obrigatório}} -- o sistema trava se estiver vazio.
        Localizado abaixo de \textit{``DADOS DO INTERMEDIÁRIO DE SERVIÇO''} e acima de \textit{``Atividade do Município''}.
    }%
}

\vspace{3pt}
\colorbox{lightGray}{%
    \parbox{\dimexpr\textwidth-2\fboxsep}{%
        \small\textbf{Texto a inserir:} \textit{Prestação de serviços odontológicos referentes a consulta e procedimentos clínicos.}
    }%
}

% --- Seção 2: Guia Completo (em duas colunas) ---
\section{Guia Completo de Preenchimento}

\begin{multicols}{2}

\subsection{Dados do Serviço e Localização}
\begin{itemize}
    \item \textbf{Atividade:} 412 -- 04.12 -- Odontologia
    \item \textbf{NBS:} 1.2301.23.00 -- Serv. odontológicos
    \item \textbf{País:} Brasil \,|\, \textbf{UF:} Goiás \,|\, \textbf{Cidade:} Goiânia
\end{itemize}

\subsection{ISSQN}
\begin{itemize}
    \item \textbf{Simples Nacional:} Optante -- ME/EPP
    \item \textbf{Regime:} Apuração tributos federais e municipais pelo Simples Nacional
    \item \textbf{Tributação:} Operação Tributável
    \item \textbf{Retenção ISSQN:} Não Retido
\end{itemize}

\subsection{Tributação Federal}
\begin{itemize}
    \item \textbf{PIS/COFINS:} 08 -- Sem Incidência
    \item \textbf{Retenção PIS/COFINS/CSLL:} Não Retidos
    \item \textbf{\textcolor{accentColor}{Perc. Tot. Tributos -- SN: 13,45}}
\end{itemize}
\colorbox{warningOrange!15}{%
    \parbox{\dimexpr\linewidth-2\fboxsep}{%
        \scriptsize\textbf{Atenção:} preencher o campo 13,45\% para evitar erro de validação.
    }%
}

\columnbreak

\subsection{IBS/CBS}
\begin{itemize}
    \item \textbf{Indicador:} 030101 -- Serviço prestado fisicamente sobre pessoa ou bem móvel em localidade determinada
    \item \textbf{Sit. Tributária:} 200 -- Alíquota reduzida
    \item \textbf{Classif. Tributária:} 200029 -- Serviços de saúde (art. 9\textsuperscript{o}, LC 196/2024)
    \item \textbf{Tipo Ente Gov.:} \textit{(vazio)}
    \item \textbf{Tipo Operação:} Fornecimento com pagamento já realizado
\end{itemize}

\subsection{Informações Complementares}
\colorbox{lightGray}{%
    \parbox{\dimexpr\linewidth-2\fboxsep}{%
        \scriptsize\texttt{DOCUMENTO EMITIDO POR ME OU EPP OPTANTE PELO SIMPLES NACIONAL; NÃO GERA DIREITO A CRÉDITO FISCAL DE IPI.}
    }%
}

\end{multicols}

% --- Seção 3: Resumo para Desenvolvedor ---
\section{Resumo para o Desenvolvedor (Documentação API)}

\vspace{2pt}
\footnotesize
\begin{tabularx}{\textwidth}{>{\columncolor{tableHeader}}l X}
    \hline
    \rowcolor{primaryColor!15}
    \textbf{Campo API} & \textbf{Valor / Regra} \\
    \hline
    \texttt{xServ} & Enviar a string completa da prestação de serviço \\
    \texttt{vTotTrib} (Lei 12.741) & Valor calculado sobre \textbf{13,45\%} \\
    \texttt{CST\_IBS\_CBS} & Código \textbf{200} com classificação \textbf{200029} \\
    Retenções (PIS, COFINS, CSLL, IRRF) & Enviar como \textbf{0,00} em formato \texttt{Double/Decimal} \\
    \hline
\end{tabularx}

\vspace{4pt}
\colorbox{successGreen!12}{%
    \parbox{\dimexpr\textwidth-2\fboxsep}{%
        \small\textbf{\textcolor{successGreen}{Dica:}} Se não encontrar o campo de descrição, ele pode estar recolhido. Clique no título da seção \textit{IDENTIFICAÇÃO DOS SERVIÇOS} para expandir.
    }%
}

\end{document}
